\subsection{Wild equivariant arcs and the exact characteristic-three range}

\begin{lemma}[A universal wild trace invariant]
\label{lem:B-RD26-trace-arc}
\status{Proved; independent complete proof and primary-source audit by the parent passed, with an independent exact certificate rerun.}
Let $n\geq1$, and let $A\geq1$ be divisible by three. Put
\[
 d=2\cdot3^{n-1}\left(A+3^n+
        \sum_{j=1}^{n-1}2\cdot3^{n-j-1}\min(A,3^{2j})\right).
                                                        \tag{B-RD26.1}
\]
There is a good supersingular elliptic curve over a complete DVR
$R$ with algebraically closed residue field of characteristic three,
absolute ramification index $d$, and full rational $3^n$-torsion,
whose primitive torsion parameters have extrema $A,1$.
Moreover $R$ and this marked curve admit faithful semilinear actions
of the following groups, fixing every marked torsion point:
\[
 C_3;\qquad C_6\ \text{if }A\text{ is odd};\qquad
 C_3\rtimes C_4\ \text{if }A\equiv1\pmod4.
\]
The action on the special fibre is the usual faithful subgroup of
$\operatorname{Aut}(Y^2=X^3-X)$, and all these actions are trivial
on the residue field. The order-four generator acts on $C_3$ by
inversion.
\end{lemma}

\begin{proof}
Let $k$ be an algebraic closure of $\mathbf F_3$, and work over
$W=W(k)$. The elliptic curve
\[
       \mathcal E_u:\quad Y^2=X^3+uX^2-X
                                                        \tag{B-RD26.2}
\]
over $W[[u]]$ is a universal framed deformation of its special
fibre $E_0$. Here the framing identifies the special fibre with
$E_0$, and isomorphisms respect that identification. To check the
parameter directly, complete the square in any deformation, move
the simple root lifting $X=0$ to zero, and take the unique fourth
root reducing to one to make the coefficient of $X$ equal to $-1$.
The only remaining coefficient is $u$. The simple-root and
fourth-root operations are unique because two and four are units.
This also proves the required uniqueness for coordinate changes
reducing to the identity. Its discriminant is $16(u^2+4)$, a unit.
Serre--Tate identifies this deformation with that of the height-two
formal group; the precise nilpotent-base theorem is
\cite[Theorem~1.2.1]{Katz1981SerreTate}.

Use the odd formal parameter $T=X^{-1/2}$ at infinity, with the
choice of sign for which
\[
 Y=-T^{-3}(1+uT^2-T^4)^{1/2},\qquad
 \omega=\frac{dX}{2Y}=
                  \frac{dT}{(1+uT^2-T^4)^{1/2}}.
\]
The coefficient of $T^3$ in $[3](T)$ is exactly $4u$.
Indeed the formal logarithm begins $T-uT^3/6$, so substitution in
$\log([3](T))=3\log(T)$ gives this coefficient. This identity is
an identity of integral coefficients, although it was calculated
after inverting three.

Let $R_n$ be the full Drinfeld-level deformation ring, and let
$x,y$ be the parameters of a universal basis of $3^n$-torsion.
It is a complete regular local ring with these two regular
parameters, by \cite[Theorem~2.1.2 and Proposition~2.2.3(i)]
{CarterStrauch2021LubinTate}. The universal family in (B-RD26.2)
pulls back to this ring. Changing the special-fibre framing gives
an action on $R_n$ and its family of
\[
 G=\langle\sigma,\tau\rangle=C_3\rtimes C_4,
\]
where on $E_0$ the first generator translates $X$ by one and the
second is the order-four automorphism with $T\mapsto iT$,
$i^2=-1$. This is the functorial action by precomposition of the
framing; its isomorphisms preserve the universal level structure.
In the explicit family, the second action is $u\mapsto-u$ and
$T\mapsto iT$ exactly. The first action reduces to
\[
              h(T)=\frac{T}{(1+T^2)^{1/2}}.
\]
The isomorphisms preserve the origin. Consequently, for $j=0,1,2$,
the difference between the series for $\sigma^j(T)$ and
\[
                 h_j(T)=\frac{T}{(1+jT^2)^{1/2}}
\]
belongs to $(3,u)T W[[u,T]]$. The factor $T$ here is essential:
it follows from preservation of the zero section, in addition to
agreement of the special-fibre series.

We record two bounds in $R_n$ needed to make the forthcoming
rational expressions integral on a chart. Set
\[
               M=3^{2n-2},\qquad H=6M.
\]
Every three-torsion coordinate belongs to
$\mathfrak m_{R_n}^{M}$, by iterating the prime factorization in
\cite[Proposition~2.2.3(vii)]{CarterStrauch2021LubinTate}.
Weierstrass preparation factors $[3](T)$ as a unit times
$T\prod_{0\ne Q\in E[3]}(T-t(Q))$. Its coefficient of $T^3$
is a sum of terms containing at least six of these coordinates.
The formula $[3]_{T^3}=4u$ therefore proves
\[
                    u\in\mathfrak m_{R_n}^{H}.          \tag{B-RD26.3}
\]
Likewise $3\in\mathfrak m_{R_n}^{8M}$. In particular $H\geq6$.
The special formal series is $[3](T)=-T^9$. For an explicit
algebraic check, the three-division polynomial of (B-RD26.2) is
$3X^4+4uX^3-6X^2-1$, and its tripling numerator reduces to $X^9$
modulo $(3,u)$. Substitution $X=T^{-2}$ with the square-root
branch whose leading coefficient is three gives the asserted
series. The sign is also checked directly by the characteristic-three
addition law: it gives $Y([3]P)=-Y^9$, using the polynomial identity
\[
 ((X^3-X)^2-1)(X-X^9)-(X^3-X)+(X^3-X)^5=0
                                      \quad\text{in }\mathbf F_3[X].
\]
Hence the successive coordinates $[3^j]y$ have weighted
initial forms $(-1)^jY^{3^{2j}}$ whenever $j<n$: the error ideal
$(3,u)$ has weight at least $H>M$, while these weights are at
most $M$.

Here is a precise weighted chart construction, including the
possible coefficients divisible by three in a Cohen presentation.
Give $x,y,3$ weights $A,1,d$. The factorization argument in the
proof of Theorem~\ref{thm:B-RD12}, which does not use tame descent
at this stage, gives lifts of all primitive torsion coordinates
with the prescribed weighted initial forms: a nonzero multiple
of $Y$, or the lower-weight term of $X+cY^{3^{2j}}$, or that
binomial at an equality of weights. Choose $c\in k^\times$ so
none of these finitely many initial forms vanishes at $(c,1)$.
In fact one may choose $c=i$ uniformly. All nonzero scalar
coefficients of the boundary binomials are in $\mathbf F_3^\times$
by the preceding explicit series and the linear terms of formal
addition and multiplication by units. The only excluded values
at a boundary are therefore $1,-1$, whereas $i^2=-1$ makes
$i\ne\pm1$. This optional choice fixes $\lambda=-1$ below.
The relation between three and the product of all primitive
coordinates supplies a generator of the Cohen kernel
\[
          R_n\simeq W[[X_0,Y_0]]/(3-P(X_0,Y_0)).        \tag{B-RD26.4}
\]
The series $P$ is a unit times the product of the chosen lifts.
It has weight at least $d$, and its weight-$d$ initial form,
evaluated at $(c,1)$, is a nonzero element $v_0\in k$.
Terms of a lifted factor divisible by three have weight at least
$d$, strictly greater than the weight of that factor. In the
whole product they consequently have weight strictly greater
than $d$. Thus the weight-$d$ part is independent of the variable
used below for $3/\pi^d$. This explicitly deals with low ordinary
degree monomials whose coefficients are divisible by three.
The generator in (B-RD26.4) has coefficient one at its linear
term $3$, so it generates the Cohen kernel, rather than merely
an ideal contained in it.

Adjoin the chart variables
\[
          x=\pi^A Z,\qquad y=\pi,\qquad3=\pi^d V,
                                                        \tag{B-RD26.5}
\]
where $Z,V$ are units, and complete at $Z=c,V=v_0,\pi=0$.
This chart has the two relations
\[
           3=\pi^d V,\qquad
           V=P(\pi^A Z,\pi)/\pi^d.                    \tag{B-RD26.6}
\]
The division in the second relation is legitimate: expand each
coefficient of $P$ in Teichmuller digits and replace each factor
$3^b$ by $\pi^{db}V^b$ using the first relation. The preceding
weighted bound gives a convergent series with nonnegative powers
of $\pi$. Its reduction is a polynomial in $Z$ alone, with value
$v_0$ at $c$. The second equation therefore has derivative one
with respect to $V$ modulo the maximal ideal. Equivalently, this
is the completed chart obtained by adjoining $x/y^A$ and $3/y^d$
to $R_n$ and removing the exceptional torsion. Before imposing
the next equation it is a complete regular local ring of
dimension two, with parameters $\pi,Z-c$. Formula (B-RD26.3)
gives $u\in\pi^H$ on this chart, while $3\in\pi^d$ and
$d\geq8M>H$.

The $C_3$ action already extends to this chart, before any trace
is formed. Indeed $\sigma x/x$ and $\sigma y/y$ are units of
$R_n$, and
\[
 \sigma Z=(\sigma x/x)(\sigma y/y)^{-A}Z,\qquad
 \sigma V=(\sigma y/y)^{-d}V.
\]
Both displayed multipliers reduce to one. The action therefore
preserves the centre $Z=c,V=v_0$ and extends continuously to its
completion. These formulas come from the actual ratios in the
fraction field, so they respect both chart relations.

In its fraction ring put
\[
 z_x=x^{-2},\quad z_y=y^{-2},\quad
 D_y=\operatorname{Tr}_{\langle\sigma\rangle}(z_y^2),
 \quad
 I=\frac{\operatorname{Tr}_{\langle\sigma\rangle}
                      (z_y^2(z_x-z_y))}{D_y},\quad
 U_y=N_{\langle\sigma\rangle}(y)^{-2}.                \tag{B-RD26.7}
\]
These are exact $C_3$-invariants. We verify all divisions and
error estimates, interpreting $O(\pi^b)$ inside the fraction
ring of this chart. The origin-preserving series above give
\[
 (\sigma^j x)^{-2}=x^{-2}+j+O(\pi^{H-2A}),\qquad
 (\sigma^j y)^{-2}=y^{-2}+j+O(\pi^{H-2}).
\]
It follows that
\[
 \begin{split}
 D_y&=3y^{-4}+6y^{-2}+5+O(\pi^{H-4})
                         =-1+O(\pi^{H-4}),\\
 I&=z_x-z_y+O(\pi^{H-2A-4}),\\
 U_y&=y^{-6}-y^{-2}+O(\pi^{H-6}).
 \end{split}                                           \tag{B-RD26.8}
\]
For the middle equality, subtract $(z_x-z_y)D_y$ before
estimating the numerator: each remaining factor has valuation
at least $H-2A-4$. For the last equality use
$\prod_{j=0}^2(y^{-2}+j)=y^{-6}+3y^{-4}+2y^{-2}$;
its difference from $y^{-6}-y^{-2}$ has a factor three.
Since $H>4$, $D_y$ is an actual unit of the chart. These
calculations also show that all the normalized expressions used
next are convergent integral series, not just valuation
estimates after an unspecified specialization.

Put $a=A/3$, and take $\lambda\in W^\times$ with residue $c^{-2}$.
Impose the exact invariant equation $I=\lambda U_y^a$.
After multiplication by $x^2$ this is an integral equation on
the chart, and (B-RD26.8) gives its explicit form
\[
 x^2(I-\lambda U_y^a)
   =1-\pi^{2A-2}Z^2-\lambda Z^2(1-\pi^4)^a
                                  +O(\pi^{H-4}).       \tag{B-RD26.9}
\]
Its reduction is $1-\bar\lambda Z^2$. At $Z=c$ its derivative
is $-2\bar\lambda c\ne0$. The formal implicit-function theorem,
applied together with the second equation of (B-RD26.6), now
eliminates $Z$ and $V$. The resulting ring $R$ is a complete DVR
with uniformizer $\pi$ and $3=\pi^d$ times a unit. This is a
simple-root argument over a complete local ring; no division by
three or averaging by the wild group order was used.
The choice of $c$ guarantees that every primitive coordinate has
exactly its assigned valuation, so the extrema are $A,1$.

Finally the chart and this branch are stable under the asserted
groups. Since every coordinate isomorphism preserves the
origin, $\sigma x/x$ and $\sigma y/y$ are units. They are both
one modulo $\pi$; hence $\sigma$ fixes the residue of $Z=x/y^A$
and of $V=3/y^d$. The rational equation is exactly
$C_3$-invariant, so its unique simple branch is stable.
For $\tau$, $x,y$ have eigenvalue $i$ exactly. Conjugation
permutes the three factors in the trace and norm, giving
\[
             \tau I=-I,\qquad \tau U_y=-U_y.
\]
If $A\equiv1\pmod4$, then $a=A/3$ is odd, the equation is
stable, and $\tau(Z)\equiv i^{1-A}c=c$. Also $4\mid d$ in
this case. For the order-two element $\tau^2$, stability of the
equation is automatic; its branch is stable when $A$ is odd,
and then $2\mid d$. The wild action is faithful on $R$:
\[
                 \sigma\pi=\pi-\tfrac12\pi^3
                                      +O(\pi^5),
\]
because the deformation error has order at least $H+1>3$.
The tame action has the stated faithful eigenvalue on $\pi$.
All actions fix $W$ and the residue field. The generic marked
elliptic curve and its level structure descend over
$\operatorname{Frac}(R)^G$ (or the indicated subgroup's fixed
field). This last assertion concerns the generic fibre; the
integral good model over $R$ is used for its faithful special
inertia action. This proves every assertion of the lemma.
\end{proof}

\begin{remark}[Exact finite certificate]
The accompanying program \path{code/verify_wild_trace_arcs.py}
checks the trace and norm identities, the special tripling
identity, the logarithmic coefficient $[3]_{T^3}=4u$, the
equivariance of the residual curve, and the exact integer
construction and least-index formula over a finite range.
It checks $66$ symbolic equivariance identities, $109727$ constructed
integer arcs, and $25000$ least-index comparisons through $e=50000$
and $n=8$. The final sign-expanded author run passed in
$2.5154165040003136$ seconds; the independent parent and
obstructions/crystals audits passed in $2.827041984997777$ and
$2.6579352160006238$ seconds respectively. Both independent
audits also checked the complete proof and original sources.
It does not computationally instantiate the universal formal
deformation or its descent; those are proved above.
\end{remark}

\begin{theorem}[B-RD26: the sharp wild higher-torsion criterion]
\label{thm:B-RD26-wild}
\status{Proved; independent parent proof, source, and exact arithmetic audits passed.}
Let $p=3$, $e$ be a positive even integer, $s=v_3(e/2)$, and
$m\in\{3,6,12\}$. An additive elliptic curve with full rational
$3^n$-torsion and exact semistability defect $m$ exists over a
complete DVR of absolute ramification index $e$ and perfect
residue field if and only if
\[
 s\geq n-1,\qquad
 e\geq\frac8m\,3^{\max(2n-2,\,4n-4-s)}.              \tag{B-RD26.10}
\]
The fields constructed here have algebraically closed residue
field. No finite-residue-field assertion is needed for this
existence criterion.
\end{theorem}

\begin{proof}
Necessity is Proposition~\ref{prop:B-RD20}. Suppose the two
inequalities hold. Start at $r_0=\max(0,2n-2-s)$ and increase
$r$ until $r=n-1$ or the inequality
\[
                 me\geq8\cdot3^{2n+r-2}
\]
would fail on increasing $r$ by one. Define
\[
       A=\frac{me}{2\cdot3^{2n-r-2}}-3^{2r+1},
       \qquad B=1.                                    \tag{B-RD26.11}
\]
Then $3\mid A$, $A\geq3^{2r}$, and, unless $r=n-1$,
$A<3^{2r+2}$. The divisibility by three follows from
$v_3(me/2)=s+1$ and $2n-r-2\leq s$.
The geometric-series identity (B-RD20.5) therefore says that
(B-RD26.1) equals $me$. For $m=6$, $A$ is odd; for $m=12$,
$A\equiv1\pmod4$. Apply Lemma~\ref{lem:B-RD26-trace-arc}
and the subgroup of order $m$. Its fixed field has absolute
index $me/m=e$. Galois descent is effective for the elliptic
curve, for instance using the invariant divisor $3[0]$.
The marked torsion points descend and are all rational.
The special-fibre action is faithful of order $m$, so the
good-reduction criterion gives exact semistability defect $m$.
In particular reduction is additive and potentially good.
\end{proof}

\begin{corollary}[B-RD27: the complete higher elliptic range]
\label{cor:B-RD27-exact-three}
\status{Proved; the full wild construction and integer criterion have passed independent audit.}
For $p=3$, write $e=2u$, $s=v_3(u)$ and
$L=\lfloor\log_3u\rfloor$. The largest level admitting an additive
elliptic curve with full rational $3^n$-torsion is
\[
 N(e)=1+\min\left(s,\left\lfloor\frac{L+1}{2}\right\rfloor,
                         \left\lfloor\frac{L+s+1}{4}\right\rfloor\right).
                                                        \tag{B-RD27.1}
\]
Equivalently, existence at level $n$ is exactly
\[
  v_3(e/2)\geq n-1,\qquad
  e\geq\frac23\,3^{\max(2n-2,4n-4-v_3(e/2))}.          \tag{B-RD27.2}
\]
Among potentially good elliptic curves, the least positive full
torsion level forcing good reduction is $N(e)+1$.
In particular full $27$-torsion with additive reduction exists
at $e=216$, with exact defect twelve. Together with the already
proved tame criterion, these formulas decide the elliptic
higher-torsion existence question at every odd prime.
\end{corollary}

\begin{proof}
Defect twelve gives the weakest wild inequality and is always
available. Its realized maximum is the upper bound $N_{\mathrm w}$
already calculated in Corollary~\ref{cor:B-RD21}. The tame maximum
is no larger. The final index statement uses the potentially good
hypothesis; multiplicative Tate examples retain the unrestricted
qualification in Proposition~\ref{prop:B-RD15}.
\end{proof}
